\section{Systemy wyszukiwania informacji}

\subsection{Ewolucja od wyszukiwania boolowskiego do semantycznego}

Tradycyjne systemy wyszukiwania informacji opierały się na modelu boolowskim, w którym wynik wyszukiwania ma charakter binarny, bez uwzględnienia stopnia dopasowania dokumentu do intencji użytkownika. Rozwój systemów typu \textit{free text retrieval} wprowadził miary takie jak TF-IDF (ang. \textit{Term Frequency-Inverse Document Frequency}), umożliwiające kwantyfikację istotności terminów w kontekście kolekcji dokumentów \cite{sparckjones1972statistical, salton1988term}. Współczynnik TF-IDF przypisuje wyższą wagę terminom występującym często w danym dokumencie, ale rzadko w całej kolekcji, co pozwala na identyfikację słów kluczowych charakterystycznych dla treści dokumentu.

Klasyczne metody reprezentacji tekstu, takie jak model \textit{Bag of Words} (BoW) czy TF-IDF, traktują jednak dokument jako nieuporządkowany zbiór słów, ignorując ich kolejność oraz zależności kontekstowe \cite{harris1954distributional}. Podejście to, określane jako reprezentacja leksykalna, nie uwzględnia podobieństwa semantycznego między wyrażeniami -- terminy ,,zawał serca'' i ,,myocardial infarction'' są traktowane jako całkowicie odrębne jednostki, mimo identycznego znaczenia.

\subsection{Modele embeddingowe}

Przełomem w kierunku reprezentacji semantycznej było wprowadzenie modelu Word2Vec \cite{mikolov2013efficient, mikolov2013distributed}, który mapuje słowa do gęstej przestrzeni wektorowej o niższym wymiarze (typowo 100--300 wymiarów). W przestrzeni tej słowa o podobnym znaczeniu lub występujące w podobnych kontekstach są reprezentowane przez wektory bliskie sobie według miary odległości, takiej jak podobieństwo kosinusowe. Model Word2Vec opiera się na hipotezie dystrybucyjnej Harrisa \cite{harris1954distributional}, zgodnie z którą znaczenie słowa jest determinowane przez kontekst jego występowania.

Ograniczeniem modelu Word2Vec jest generowanie statycznych embeddingów -- każde słowo posiada pojedynczą reprezentację wektorową niezależnie od kontekstu użycia. Problem ten rozwiązują kontekstowe modele językowe oparte na architekturze Transformer \cite{vaswani2017attention}. Model BERT (ang. \textit{Bidirectional Encoder Representations from Transformers}) \cite{devlin2019bert} generuje reprezentacje słów uwzględniające dwukierunkowy kontekst całego zdania, co pozwala na rozróżnienie znaczeń polisemicznych. Adaptacja modelu BERT do zadania generowania embeddingów zdaniowych zaowocowała powstaniem architektury Sentence-BERT \cite{reimers2019sentence}, umożliwiającej efektywne porównywanie podobieństwa semantycznego całych sekwencji tekstowych.

Współczesne badania zmierzają w kierunku multimodalnych przestrzeni embeddingowych, łączących reprezentacje różnych typów danych. Model CLIP \cite{radford2021learning} tworzy wspólną przestrzeń wektorową dla tekstu i obrazów, umożliwiając wyszukiwanie obrazów na podstawie zapytań tekstowych i odwrotnie. Podobne podejście reprezentują modele takie jak ALIGN \cite{jia2021scaling} czy współczesne architektury embeddingowe, takie jak Jina Embeddings \cite{gunther2023jina}, oferujące zunifikowaną przestrzeń dla wielu modalności i języków.

Embedding definiowany jest formalnie jako funkcja $f: \mathcal{X} \rightarrow \mathbb{R}^d$, która mapuje obiekty ze zbioru wejściowego $\mathcal{X}$ (słowa, zdania, dokumenty, obrazy) do $d$-wymiarowej przestrzeni wektorowej, przy czym podobieństwo semantyczne obiektów jest zachowane jako bliskość geometryczna ich reprezentacji wektorowych. Miarą podobieństwa najczęściej stosowaną w praktyce jest podobieństwo kosinusowe:
\begin{equation}
    \text{sim}(\mathbf{u}, \mathbf{v}) = \frac{\mathbf{u} \cdot \mathbf{v}}{\|\mathbf{u}\| \|\mathbf{v}\|}
\end{equation}
gdzie $\mathbf{u}$ i $\mathbf{v}$ są wektorami embeddingów porównywanych obiektów.

\subsection{Architektura RAG}

Duże modele językowe (ang. \textit{Large Language Models}, LLM), takie jak GPT \cite{brown2020language} czy LLaMA \cite{touvron2023llama}, pomimo imponujących zdolności generacyjnych, charakteryzują się istotnymi ograniczeniami w kontekście zastosowań specjalistycznych. Wiedza tych modeli jest zamrożona w momencie zakończenia treningu, co uniemożliwia uwzględnienie informacji aktualnych lub wysoce specjalistycznych. Ponadto modele te są podatne na zjawisko halucynacji -- generowania treści brzmiących wiarygodnie, lecz niezgodnych z faktami \cite{ji2023survey}.

Architektura RAG (ang. \textit{Retrieval-Augmented Generation}) \cite{lewis2020retrieval} stanowi rozwiązanie tych problemów poprzez integrację komponentu wyszukiwania informacji z modelem generacyjnym. System RAG składa się z trzech głównych faz: indeksowania, wyszukiwania oraz generacji. Schemat architektury przedstawiono na Rysunku \ref{fig:ragArchitecture}.

W fazie indeksowania dokumenty źródłowe są przetwarzane i dzielone na fragmenty (ang. \textit{chunks}) o określonej długości. Każdy fragment jest następnie przekształcany przez model embeddingowy w reprezentację wektorową i zapisywany w wektorowej bazie danych wraz z metadanymi umożliwiającymi identyfikację źródła \cite{johnson2019billion}.

Faza wyszukiwania inicjowana jest przez zapytanie użytkownika, które jest transformowane do tej samej przestrzeni wektorowej co zindeksowane dokumenty. Na podstawie miary podobieństwa (typowo podobieństwa kosinusowego lub odległości euklidesowej) system identyfikuje $k$ najbardziej podobnych fragmentów dokumentów.

W fazie generacji model językowy otrzymuje kontekst składający się z wyszukanych fragmentów dokumentów oraz oryginalnego zapytania użytkownika. Na tej podstawie generuje odpowiedź, która może być weryfikowana przez użytkownika poprzez porównanie z dostarczonymi źródłami.

Architektura RAG oferuje szereg zalet w kontekście zastosowań specjalistycznych. Po pierwsze, umożliwia aktualizację bazy wiedzy bez konieczności ponownego trenowania modelu językowego -- wystarczy dodanie nowych dokumentów do indeksu wektorowego. Po drugie, zapewnia wyjaśnialność poprzez możliwość wskazania konkretnych fragmentów dokumentów wykorzystanych do wygenerowania odpowiedzi. Po trzecie, redukuje ryzyko halucynacji poprzez zakotwiczenie generowanej treści w dostarczonym kontekście \cite{shuster2021retrieval}.

\begin{figure}[htbp]
    \centering
    \begin{tikzpicture}[
    scale=0.65,
    transform shape,
    font=\small,
    block/.style={
        draw,
        rounded corners=2pt,
        thick,
        minimum width=2.8cm,
        minimum height=1.0cm,
        align=center
    },
    data/.style={block, fill=gray!12},
    embm/.style={block, fill=purple!20},
    store/.style={block, fill=orange!25},
    ann/.style={block, fill=yellow!25},
    gen/.style={block, fill=green!20},
    ctx/.style={block, fill=teal!25},
    arrow/.style={-{Stealth[]}, thick},
    phase/.style={draw=gray, dashed, rounded corners=4pt, inner sep=0.4cm}
]

% ============================================================
% FAZA INDEKSOWANIA
% ============================================================
\node[data] (docs) {Dokumenty\\źródłowe};
\node[data, right=2.0cm of docs] (chunks) {Podział na\\fragmenty};
\node[embm, right=2.0cm of chunks] (demb) {Model\\embeddingowy};
\node[store, right=2.0cm of demb] (db) {Wektorowa\\baza danych};

\draw[arrow] (docs) -- (chunks);
\draw[arrow] (chunks) -- (demb);
\draw[arrow] (demb) -- node[below, font=\scriptsize, text=gray] {$1\times d_{\mathrm{emb}}$} (db);

% ============================================================
% FAZA WYSZUKIWANIA
% ============================================================
\node[data, below=2.6cm of docs] (query) {Zapytanie\\użytkownika};
\node[embm, right=2.0cm of query] (qemb) {Model\\embeddingowy};
\node[data, right=2.0cm of qemb] (qvec) {Wektor\\zapytania};
\node[ann, right=2.0cm of qvec] (search) {Wyszukiwanie ANN};

\draw[arrow] (query) -- (qemb);
\draw[arrow] (qemb) -- node[below, font=\scriptsize, text=gray] {$1\times d_{\mathrm{emb}}$} (qvec);
\draw[arrow] (qvec) -- (search);
\draw[arrow] (db) -- node[right, font=\scriptsize] {indeks} (search);

% ============================================================
% FAZA GENERACJI
% ============================================================
\node[data, below=2.6cm of query] (topk) {Top-$k$\\fragmentów};
\node[ctx, right=2.0cm of topk] (aug) {Augmentacja\\kontekstu};
\node[gen, right=2.0cm of aug] (llm) {LLM};
\node[data, right=2.0cm of llm] (ans) {Wygenerowana\\odpowiedź};

\draw[arrow] (search.south) -- ++(0,-0.8) -| (topk.north);
\draw[arrow] (topk) -- (aug);
\draw[arrow] (aug) -- (llm);
\draw[arrow] (llm) -- (ans);

% zapytanie użytkownika trafia również do kontekstu promptu
\draw[arrow, gray] (query.south east) -- node[sloped, above, font=\scriptsize, text=gray] {zapytanie} (aug.north west);

% ============================================================
% RAMKI FAZ
% ============================================================
\node[phase, fit=(docs)(db),
      label={[font=\small\itshape, gray]above:Faza indeksowania}] {};
\node[phase, fit=(query)(search),
      label={[font=\small\itshape, gray]above left:Faza wyszukiwania}] {};
\node[phase, fit=(topk)(ans),
      label={[font=\small\itshape, gray]below left:Faza generacji}] {};

\end{tikzpicture}

    \caption{Architektura RAG (\textit{Retrieval-Augmented Generation}) z podziałem na fazy indeksowania, wyszukiwania i generacji. Zapytanie użytkownika trafia zarówno do modelu embeddingowego, jak i -- wraz z wyszukanymi fragmentami -- do kontekstu modelu językowego.}
    \label{fig:ragArchitecture}
\end{figure}


% ---------------------------------------------------
\section{Modele automatycznego rozpoznawania mowy}
\label{sec:modeleASR}

\subsection{Historia i rozwój modeli ASR}

Automatyczne rozpoznawanie mowy (ang. \textit{Automatic Speech Recognition}, ASR) stanowi jedno z klasycznych zagadnień przetwarzania języka naturalnego, którego historia sięga lat siedemdziesiątych XX wieku. Wczesne systemy ASR opierały się na ukrytych modelach Markowa (ang. \textit{Hidden Markov Models}, HMM), które modelowały sygnał mowy jako sekwencję dyskretnych stanów akustycznych \cite{rabiner1989tutorial}. W podejściu tym zakłada się, że obserwowany sygnał akustyczny jest generowany przez ukryty proces stochastyczny, którego stany odpowiadają jednostkom fonetycznym.

Rozszerzeniem modeli HMM były systemy hybrydowe GMM-HMM, łączące ukryte modele Markowa z mieszaninami gaussowskimi (ang. \textit{Gaussian Mixture Models}) \cite{gales2008application}. Mieszaniny gaussowskie służyły do modelowania rozkładu prawdopodobieństwa emisji w poszczególnych stanach HMM, co pozwalało na lepsze uchwycenie zmienności akustycznej sygnału mowy. Systemy te wymagały jednak znacznej wiedzy eksperckiej. Konieczne było ręczne projektowanie cech akustycznych (np. współczynników MFCC), tworzenie słowników wymowy oraz definiowanie modeli fonologicznych dla danego języka.

Przełom w dziedzinie ASR nastąpił wraz z zastosowaniem głębokich sieci neuronowych. Początkowo sieci te zastępowały komponent GMM w systemach hybrydowych DNN-HMM \cite{hinton2012deep}, zachowując strukturę modelowania sekwencji opartą na HMM. Wprowadzenie mechanizmu uwagi (ang. \textit{attention mechanism}) \cite{bahdanau2014neural} oraz architektur typu sekwencja-do-sekwencji (ang. \textit{sequence-to-sequence}) \cite{sutskever2014sequence} umożliwiło bezpośrednie modelowanie zależności między sekwencją akustyczną a sekwencją tekstową bez konieczności jawnego definiowania wyrównania czasowego.

Architektura Transformer \cite{vaswani2017attention}, oparta wyłącznie na mechanizmie uwagi bez komponentów rekurencyjnych, znalazła szerokie zastosowanie w modelach ASR nowej generacji. Modele te, określane jako end-to-end, eliminują tradycyjny pipeline przetwarzania składający się z oddzielnych modułów akustycznego, leksykalnego i językowego. Zamiast tego realizują bezpośrednie mapowanie surowego sygnału audio na sekwencję tekstową, co upraszcza architekturę systemu i eliminuje propagację błędów między modułami.

\subsection{Architektura współczesnych modeli ASR}

Współczesne modele ASR można podzielić na dwie główne kategorie ze względu na metodę treningu: modele wykorzystujące uczenie samonadzorowane (ang. \textit{self-supervised learning}) oraz modele trenowane w sposób nadzorowany (ang. \textit{supervised learning}) na dużych zbiorach transkrybowanych nagrań.

\subsubsection{Wav2Vec 2.0}

Model Wav2Vec 2.0 \cite{baevski2020wav2vec} reprezentuje podejście oparte na uczeniu samonadzorowanym. Architektura składa się z trzech głównych komponentów: enkodera konwolucyjnego (ang. \textit{Convolutional Neural Network}, CNN), transformera oraz modułu kwantyzacji. Enkoder konwolucyjny przetwarza surowy sygnał audio (próbkowany typowo z częstotliwością 16 kHz) i przekształca go w sekwencję wektorów cech o niższej rozdzielczości czasowej. Następnie transformer przetwarza tę sekwencję, generując kontekstowe reprezentacje uwzględniające zależności długozasięgowe. Schemat architektury przedstawiono na Rysunku \ref{fig:wav2vecArchitecture}.

Proces uczenia samonadzorowanego w modelu Wav2Vec 2.0 opiera się na maskowaniu fragmentów sekwencji wejściowej oraz uczeniu kontrastywnym (ang. \textit{contrastive learning}). Model jest trenowany do identyfikacji prawidłowych kwantyzowanych reprezentacji zamaskowanych segmentów ze zbioru negatywnych przykładów. Podejście to pozwala na wykorzystanie ogromnych ilości nieetykietowanych nagrań audio do uczenia semantycznie bogatych reprezentacji akustycznych, bez konieczności korzystania z dużych ilości transkrybowanych nagrań.

Adaptacja modelu do zadania transkrypcji wymaga etapu dostrajania (ang. \textit{fine-tuning}) z wykorzystaniem stosunkowo niewielkiego zbioru danych etykietowanych. Eksperymenty wykazały, że model Wav2Vec 2.0 osiąga konkurencyjne wyniki przy użyciu zaledwie kilkudziesięciu minut transkrybowanych nagrań \cite{baevski2020wav2vec}. Stanowi to istotną zaletę w kontekście języków i domen o ograniczonych zasobach etykietowanych. Dzięki temu podejściu Wav2Vec 2.0 znacząco zredukował zapotrzebowanie na dane z etykietami, jednocześnie poprawiając jakość rozpoznawania mowy w warunkach szumu, zmienności mówców i zróżnicowanego otoczenia akustycznego.

\begin{figure}[htbp]
    \centering
    \begin{tikzpicture}[
    scale=0.6,
    transform shape,
    font=\small,
    block/.style={
        draw,
        rounded corners=2pt,
        thick,
        minimum width=2.6cm,
        minimum height=1.0cm,
        align=center
    },
    data/.style={block, fill=gray!12},
    model/.style={block, fill=blue!20},
    trans/.style={block, fill=purple!20},
    op/.style={block, fill=orange!25},
    quant/.style={block, fill=teal!25},
    loss/.style={block, fill=yellow!25},
    arrow/.style={-{Stealth[]}, thick},
    phase/.style={draw=gray, dashed, rounded corners=4pt, inner sep=0.4cm}
]

% ============================================================
% GŁÓWNY PRZEPŁYW
% ============================================================
\node[data] (audio) {Surowy sygnał\\audio};
\node[model, right=1.0cm of audio] (cnn) {Enkoder CNN};
\node[data, right=1.0cm of cnn] (z) {Wektory cech\\$z_{1..T}$};
\node[op, right=1.0cm of z] (mask) {Maskowanie};
\node[trans, right=1.0cm of mask] (tr) {Transformer};
\node[data, right=1.0cm of tr] (c) {Reprezentacje\\kontekstowe $c_{1..T}$};

\draw[arrow] (audio) -- (cnn);
\draw[arrow] (cnn) -- (z);
\draw[arrow] (z) -- (mask);
\draw[arrow] (mask) -- (tr);
\draw[arrow] (tr) -- (c);

% ============================================================
% PRETRENING: kwantyzacja + zadanie kontrastywne
% ============================================================
\node[quant, below=2.2cm of z] (q) {Moduł kwantyzacji\\$q_{1..T}$};
\node[loss, below=2.2cm of c] (contrast) {Strata\\kontrastywna};

\draw[arrow] (z) -- (q);
\draw[arrow] (q) -- (contrast);
\draw[arrow] (c) -- (contrast);

% ============================================================
% DOSTRAJANIE: głowica klasyfikacji
% ============================================================
\node[model, above=2.2cm of c] (head) {Warstwa\\klasyfikacji (CTC)};
\node[data, right=1.0cm of head] (text) {Tekst};

\draw[arrow] (c) -- (head);
\draw[arrow] (head) -- (text);

% ============================================================
% RAMKI FAZ
% ============================================================
\node[phase, fit=(q)(contrast),
      label={[font=\small\itshape, gray]below:Pretrening (uczenie samonadzorowane)}] {};
\node[phase, fit=(head)(text),
      label={[font=\small\itshape, gray]above:Dostrajanie (\textit{fine-tuning})}] {};

\end{tikzpicture}

    \caption{Architektura modelu Wav2Vec 2.0. Podczas pretreningu zamaskowane reprezentacje kontekstowe są porównywane z kwantyzowanymi reprezentacjami cech w ramach zadania kontrastywnego. Etap dostrajania dodaje warstwę klasyfikacji generującą transkrypcję.}
    \label{fig:wav2vecArchitecture}
\end{figure}

\subsubsection{HuBERT}

Model HuBERT (ang. \textit{Hidden-Unit BERT}) \cite{hsu2021hubert} stanowi rozwinięcie koncepcji uczenia samonadzorowanego dla reprezentacji mowy. Podobnie jak Wav2Vec 2.0, HuBERT wykorzystuje architekturę składającą się z enkodera konwolucyjnego oraz transformera. Zasadnicza różnica dotyczy jednak metody uczenia.

W modelu Wav2Vec 2.0 uczenie opiera się na zadaniu kontrastywnym, w którym model rozróżnia prawdziwe reprezentacje zamaskowanych segmentów od negatywnych przykładów. HuBERT natomiast wykorzystuje podejście inspirowane modelem BERT \cite{devlin2019bert} z dziedziny przetwarzania języka naturalnego. Model jest trenowany do przewidywania dyskretnych pseudoetykiet przypisanych zamaskowanym fragmentom sygnału audio. Schemat architektury przedstawiono na Rysunku \ref{fig:hubertArchitecture}.

Pseudoetykiety generowane są w procesie iteracyjnym. W pierwszej iteracji stosuje się algorytm k-średnich (ang. \textit{k-means clustering}) do cech MFCC surowego sygnału audio, uzyskując początkowy zbiór klastrów. Model HuBERT jest następnie trenowany do przewidywania przypisań klastrowych dla zamaskowanych segmentów. W kolejnych iteracjach pseudoetykiety są aktualizowane poprzez zastosowanie algorytmu k-średnich do reprezentacji cech ukrytych wytrenowanego modelu, co prowadzi do coraz bardziej semantycznych grupowań.

Zaletą tego podejścia jest możliwość uczenia bez konieczności definiowania negatywnych przykładów, co upraszcza proces treningowy. Ponadto iteracyjne doskonalenie pseudoetykiet pozwala modelowi na stopniowe odkrywanie struktury językowej sygnału mowy. Eksperymenty wykazały, że HuBERT osiąga wyniki porównywalne lub lepsze od Wav2Vec 2.0 w zadaniach rozpoznawania mowy, szczególnie w scenariuszach z ograniczoną ilością danych etykietowanych \cite{hsu2021hubert}.

Wariantem modelu HuBERT jest mHuBERT \cite{mhubert}, przystosowany do przetwarzania wielojęzycznego, co rozszerza zastosowania oryginalnego modelu na języki o ograniczonych zasobach.

\begin{figure}[htbp]
    \centering
    \begin{tikzpicture}[
    scale=0.6,
    transform shape,
    font=\small,
    block/.style={
        draw,
        rounded corners=2pt,
        thick,
        minimum width=2.8cm,
        minimum height=1.0cm,
        align=center
    },
    data/.style={block, fill=gray!12},
    model/.style={block, fill=blue!20},
    trans/.style={block, fill=purple!20},
    op/.style={block, fill=orange!25},
    clust/.style={block, fill=teal!25},
    lab/.style={block, fill=yellow!25},
    arrow/.style={-{Stealth[]}, thick},
    phase/.style={draw=gray, dashed, rounded corners=4pt, inner sep=0.4cm}
]

% ============================================================
% UCZENIE Z MASKOWANIEM
% ============================================================
\node[data] (audio) {Sygnał\\audio};
\node[model, right=1.0cm of audio] (cnn) {Enkoder CNN};
\node[op, right=1.0cm of cnn] (mask) {Maskowanie};
\node[trans, right=1.0cm of mask] (tr) {Transformer};
\node[model, right=1.6cm of tr] (pred) {Predykcja\\pseudoetykiet};

\draw[arrow] (audio) -- (cnn);
\draw[arrow] (cnn) -- (mask);
\draw[arrow] (mask) -- (tr);
\draw[arrow] (tr) -- (pred);

% ============================================================
% GENEROWANIE PSEUDOETYKIET
% ============================================================
\node[data, above=2.6cm of mask] (mfcc) {Cechy MFCC /\\reprezentacje ukryte};
\node[clust, above=2.6cm of tr] (km) {Grupowanie\\$k$-średnich};
\node[lab, above=2.6cm of pred] (pseudo) {Pseudoetykiety};

\draw[arrow] (audio.north) |- (mfcc.west);
\draw[arrow] (mfcc) -- (km);
\draw[arrow] (km) -- (pseudo);

% kolejne iteracje: reprezentacje transformera zasilają ponowne grupowanie
\draw[arrow, dashed, gray] (tr.north) --
    node[right, font=\scriptsize, text=gray, align=left] {kolejne\\iteracje} (km.south);

% cel predykcji
\draw[{Stealth[]}-{Stealth[]}, dotted, thick, gray] (pseudo.south) --
    node[right, font=\scriptsize\itshape, text=gray] {entropia krzyżowa} (pred.north);

% ============================================================
% RAMKI FAZ
% ============================================================
\node[phase, fit=(mfcc)(pseudo),
      label={[font=\small\itshape, gray]above:Generowanie pseudoetykiet (iteracyjne)}] {};
\node[phase, fit=(audio)(pred),
      label={[font=\small\itshape, gray]below:Uczenie z maskowaniem}] {};

\end{tikzpicture}

    \caption{Architektura modelu HuBERT. Pseudoetykiety powstają przez grupowanie $k$-średnich cech MFCC (pierwsza iteracja), a w kolejnych iteracjach -- reprezentacji ukrytych wytrenowanego modelu. Model uczy się przewidywać pseudoetykiety zamaskowanych segmentów.}
    \label{fig:hubertArchitecture}
\end{figure}

\subsubsection{Whisper}

Model Whisper \cite{whisper} reprezentuje odmienne podejście, oparte na uczeniu nadzorowanym w dużej skali. Model został wytrenowany na zbiorze 680 tysięcy godzin transkrybowanych nagrań audio pozyskanych z internetu, obejmujących 99 języków. Architektura Whisper składa się z enkodera i dekodera, oba oparte na transformerze.

W przeciwieństwie do modelu Wav2Vec 2.0, Whisper operuje na spektogramach w skali melowej. Stanowią one reprezentację częstotliwościową sygnału audio uwzględniającą charakterystykę ludzkiego słuchu. Takie podejście umożliwia normalizację dynamiki i redukcję szumu na etapie ekstrakcji cech, a jednocześnie przenosi ciężar modelowania z niskopoziomowych cech akustycznych na wyższe zależności językowe i kontekstowe. Enkoder przetwarza melspektrogram i generuje sekwencję reprezentacji ukrytych. Dekoder, działający w sposób autoregresyjny, generuje kolejne tokeny tekstu na podstawie reprezentacji z enkodera oraz wcześniej wygenerowanych tokenów. Schemat architektury przedstawiono na Rysunku \ref{fig:whisperArchitecture}.

Istotną cechą modelu Whisper jest uczenie wielozadaniowe (ang. \textit{multitask learning}). Model jest trenowany równocześnie do realizacji kilku zadań: transkrypcji mowy, detekcji języka, znaczników czasowych oraz tłumaczenia mowy na język angielski. Zadanie jest specyfikowane poprzez specjalne tokeny kontrolne umieszczane na początku sekwencji wyjściowej.

Zaletą modelu Whisper jest gotowość do użycia bez konieczności dostrajania (ang. \textit{zero-shot}). Dzięki treningowi na zróżnicowanym zbiorze danych model wykazuje odporność na szum, różne warunki akustyczne oraz akcenty. Wielojęzyczność modelu pozwala na transkrypcję nagrań w wielu językach, w tym w języku polskim, bez dodatkowej adaptacji.

\begin{figure}[htbp]
    \centering
    \begin{tikzpicture}[
    scale=0.6,
    transform shape,
    font=\small,
    block/.style={
        draw,
        rounded corners=2pt,
        thick,
        minimum width=2.8cm,
        minimum height=1.0cm,
        align=center
    },
    data/.style={block, fill=gray!12},
    feat/.style={block, fill=orange!25},
    enc/.style={block, fill=blue!20},
    dec/.style={block, fill=purple!20},
    ctrl/.style={block, fill=teal!25},
    arrow/.style={-{Stealth[]}, thick},
    phase/.style={draw=gray, dashed, rounded corners=4pt, inner sep=0.4cm}
]

\node[data] (audio) {Sygnał\\audio};
\node[feat, right=1.0cm of audio] (mel) {Melspektrogram};
\node[enc, right=1.0cm of mel] (enc) {Enkoder\\(Transformer)};
\node[data, right=1.0cm of enc] (hid) {Reprezentacje\\ukryte};
\node[dec, right=2.8cm of hid] (dec) {Dekoder\\(autoregresyjny)};
\node[data, right=1.6cm of dec] (out) {Tokeny wyjściowe\\(tekst)};

\draw[arrow] (audio) -- (mel);
\draw[arrow] (mel) -- (enc);
\draw[arrow] (enc) -- (hid);
\draw[arrow] (hid) -- node[above, font=\scriptsize] {\textit{cross-attention}} (dec);
\draw[arrow] (dec) -- (out);

% tokeny kontrolne zadania
\node[ctrl, below=2.2cm of dec] (ctrl)
    {Tokeny kontrolne\\(język, zadanie, znaczniki)};
\draw[arrow] (ctrl) -- (dec);

% sprzężenie autoregresyjne
\draw[arrow, gray] (out.north) -- ++(0,1.0)
    -| node[pos=0.25, above, font=\scriptsize, text=gray] {poprzednie tokeny (autoregresja)} (dec.north);

\end{tikzpicture}

    \caption{Architektura modelu Whisper. Enkoder przetwarza melspektrogram, a autoregresyjny dekoder generuje kolejne tokeny na podstawie reprezentacji enkodera, tokenów kontrolnych określających zadanie oraz wcześniej wygenerowanych tokenów.}
    \label{fig:whisperArchitecture}
\end{figure}

\subsubsection{Data2Vec}

Model Data2Vec \cite{data2vec} stanowi uogólnienie podejść do uczenia samonadzorowanego, oferując zunifikowaną architekturę zdolną do przetwarzania różnych modalności: mowy, obrazów oraz tekstu. Kluczową innowacją jest rezygnacja z przewidywania celów specyficznych dla modalności na rzecz predykcji kontekstowych reprezentacji latentnych.

Architektura Data2Vec opiera się na mechanizmie samodestylacji z wykorzystaniem modelu nauczyciela (ang. \textit{teacher}) oraz ucznia (ang. \textit{student}). Model nauczyciela przetwarza pełne dane wejściowe i generuje reprezentacje latentne. Model ucznia, otrzymując zamaskowane dane, przewiduje reprezentacje wygenerowane przez nauczyciela dla zamaskowanych fragmentów. Wagi nauczyciela są aktualizowane jako wykładnicza średnia krocząca (EMA) wag ucznia, co zapewnia stabilność procesu uczenia. Schemat architektury przedstawiono na Rysunku \ref{fig:data2vecArchitecture}.

W przypadku mowy Data2Vec wykorzystuje architekturę transformera analogiczną do Wav2Vec 2.0. Eksperymenty na zbiorze LibriSpeech wykazały, że model osiąga niższy wskaźnik WER niż Wav2Vec 2.0 i HuBERT, szczególnie przy ograniczonej ilości danych etykietowanych \cite{data2vec}. Zunifikowane podejście do uczenia upraszcza rozwój modeli multimodalnych i umożliwia transfer wiedzy między modalnościami.

\begin{figure}[htbp]
    \centering
    \begin{tikzpicture}[
    scale=0.6,
    transform shape,
    font=\small,
    block/.style={
        draw,
        rounded corners=2pt,
        thick,
        minimum width=2.8cm,
        minimum height=1.0cm,
        align=center
    },
    data/.style={block, fill=gray!12},
    teach/.style={block, fill=blue!20},
    stud/.style={block, fill=purple!20},
    op/.style={block, fill=orange!25},
    loss/.style={block, fill=yellow!25},
    arrow/.style={-{Stealth[]}, thick},
    phase/.style={draw=gray, dashed, rounded corners=4pt, inner sep=0.4cm}
]

% ============================================================
% WSPÓLNE WEJŚCIE
% ============================================================
\node[data] (src) {Dane wejściowe\\(mowa / obraz / tekst)};

% ============================================================
% GAŁĄŹ NAUCZYCIELA
% ============================================================
\node[data, above right=1.0cm and 1.4cm of src] (full) {Pełne\\dane};
\node[teach, right=1.2cm of full] (enct) {Enkoder\\(nauczyciel)};
\node[data, right=1.2cm of enct] (rept) {Reprezentacje\\docelowe};

\draw[arrow] (src.east) -- ++(0.6,0) |- (full.west);
\draw[arrow] (full) -- (enct);
\draw[arrow] (enct) -- (rept);

% ============================================================
% GAŁĄŹ UCZNIA
% ============================================================
\node[op, below right=1.0cm and 1.4cm of src] (mask) {Maskowanie};
\node[stud, right=1.2cm of mask] (encs) {Enkoder\\(uczeń)};
\node[data, right=1.2cm of encs] (preds) {Predykcja\\reprezentacji};

\draw[arrow] (src.east) -- ++(0.6,0) |- (mask.west);
\draw[arrow] (mask) -- (encs);
\draw[arrow] (encs) -- (preds);

% ============================================================
% STRATA I AKTUALIZACJA EMA
% ============================================================
\node[loss, right=1.4cm of rept.east |- src] (loss) {Strata\\regresyjna};

\draw[arrow] (rept.east) -| (loss.north);
\draw[arrow] (preds.east) -| (loss.south);

\draw[arrow, dashed, gray] (encs.north) --
    node[right, font=\scriptsize, text=gray, align=left] {EMA\\(aktualizacja wag)} (enct.south);

% ============================================================
% RAMKI GAŁĘZI
% ============================================================
\node[phase, fit=(full)(rept),
      label={[font=\small\itshape, gray]above:Nauczyciel (\textit{teacher})}] {};
\node[phase, fit=(mask)(preds),
      label={[font=\small\itshape, gray]below:Uczeń (\textit{student})}] {};

\end{tikzpicture}

    \caption{Architektura modelu Data2Vec. Nauczyciel przetwarza pełne dane wejściowe i wyznacza kontekstowe reprezentacje latentne stanowiące cel predykcji, natomiast uczeń otrzymuje dane zamaskowane. Wagi nauczyciela są wykładniczą średnią kroczącą (EMA) wag ucznia.}
    \label{fig:data2vecArchitecture}
\end{figure}

\subsection{Adaptacja do domeny specjalistycznej}

Pomimo wysokiej jakości modeli ogólnego przeznaczenia, zastosowania w domenach specjalistycznych wymagają często dodatkowej adaptacji. Terminologia medyczna, prawnicza czy techniczna jest słabo reprezentowana w ogólnych zbiorach treningowych, co prowadzi do błędów rozpoznawania specjalistycznych terminów.

Przykładem takiej adaptacji jest projekt ADMEDVoice \cite{ADMEDVoice}, w ramach którego przeprowadzono dostrajanie modelu Whisper Medium do transkrypcji polskiej mowy medycznej. Zbiór treningowy obejmował nagrania 28 mówców w różnych scenariuszach akustycznych, zawierające terminologię kliniczną. Dostrajanie pozwoliło na redukcję wskaźnika błędu słów (ang. \textit{Word Error Rate}, WER) z 24\% do 9,5\%, co stanowi znaczącą poprawę jakości transkrypcji w domenie medycznej.

Wyniki te ilustrują znaczenie danych domenowych dla jakości systemów ASR w zastosowaniach specjalistycznych. Jednocześnie wskazują na wyzwanie związane z ograniczoną dostępnością takich zbiorów, szczególnie dla języków innych niż angielski.

Analogiczny problem dotyczy modeli embeddingowych i językowych. Rozwiązaniem jest adaptacja modeli do konkretnych domen poprzez dodatkowy trening na specjalistycznych korpusach. W dziedzinie przetwarzania języka naturalnego powstały modele takie jak BioBERT \cite{lee2020biobert}, pretrenowany na korpusach tekstów biomedycznych z PubMed, czy PubMedBERT \cite{gu2021domain}, trenowany od podstaw wyłącznie na tekstach z domeny biomedycznej. Modele te wykazują znacząco lepszą jakość w zadaniach związanych z tekstami medycznymi w porównaniu z modelami ogólnego przeznaczenia.

\subsection{Porównanie podejść}

Modele Wav2Vec 2.0, HuBERT oraz Whisper reprezentują różne podejścia do problemu ASR. Wav2Vec 2.0 i HuBERT wykorzystują uczenie samonadzorowane, co pozwala na efektywne wykorzystanie nieetykietowanych nagrań audio. Różnią się jednak metodą uczenia: Wav2Vec 2.0 stosuje uczenie kontrastywne, podczas gdy HuBERT opiera się na przewidywaniu pseudoetykiet w sposób analogiczny do modelu BERT. Whisper reprezentuje odmienne podejście, oferując wysoką jakość bez konieczności dostrajania dzięki treningowi nadzorowanemu na masywnym zbiorze danych.

Otwartym pytaniem pozostaje, czy reprezentacje ukryte generowane przez modele ASR zawierają informację semantyczną wykraczającą poza poziom akustyczny i fonetyczny. Zagadnienie to stanowi motywację dla badań nad bezpośrednim wykorzystaniem reprezentacji modeli ASR w zadaniach wyszukiwania informacji.


% ---------------------------------------------------
\section{Wyjaśnialność modeli wykorzystujących uczenie głębokie}

Modele oparte na uczeniu głębokim realizują odwzorowanie danych wejściowych na wyjściowe poprzez sekwencję nieliniowych transformacji, których pojedyncze parametry nie mają bezpośredniej interpretacji. Nurt badań określany jako wyjaśnialność (ang. \textit{explainability}) obejmuje szerokie spektrum metod, od atrybucji cech dla pojedynczej predykcji po analizę własności całych przestrzeni reprezentacji. Z perspektywy niniejszej pracy istotny jest wyłącznie ten drugi obszar. Pytanie badawcze nie dotyczy bowiem uzasadnienia konkretnej transkrypcji, lecz tego, jaka informacja jest zakodowana w reprezentacjach ukrytych enkodera modelu ASR oraz w jaki sposób można tę informację zmierzyć.

\subsection{Semantyczność reprezentacji}

Pojęcie semantycznej reprezentacji wywodzi się z hipotezy dystrybucyjnej sformułowanej przez Harrisa, zgodnie z którą znaczenie słowa wynika z kontekstów jego występowania \cite{harris1954distributional}. Hipoteza ta stanowiła podstawę modeli generujących statyczne osadzenia wektorowe, takich jak Word2Vec \cite{mikolov2013efficient, mikolov2013distributed}, w których relacje semantyczne odwzorowywane są przez zależności geometryczne między wektorami. Wprowadzenie architektury transformera \cite{vaswani2017attention} i modeli takich jak BERT \cite{devlin2019bert} pozwoliło na generowanie reprezentacji kontekstowych, a modele typu Sentence-BERT \cite{reimers2019sentence} rozszerzyły to podejście na poziom całych zdań.

Zasadniczą trudnością jest fakt, że semantyczność nie jest wielkością obserwowalną bezpośrednio. Reprezentacja wektorowa nie posiada etykiety opisującej zakodowane w niej znaczenie, wobec czego jej ocena wymaga metod pośrednich. W literaturze wyróżnić można trzy grupy takich metod: analizę geometrycznych własności przestrzeni reprezentacji, zadania probingowe wykorzystujące dodatkowy model nadzorowany oraz dekompozycję reprezentacji na interpretowalne komponenty.

\subsection{Geometria przestrzeni reprezentacji}

Analiza geometryczna opiera się na założeniu, że rozkład reprezentacji w przestrzeni wielowymiarowej odzwierciedla jakość zakodowanej w nich informacji. Ethayarajh wykazał, że reprezentacje kontekstowe modeli BERT, ELMo i GPT-2 zajmują wąski stożek przestrzeni, przez co podobieństwo kosinusowe dowolnych dwóch losowych reprezentacji przyjmuje wartości znacząco dodatnie \cite{cone}. Zjawisko to, określane jako anizotropia przestrzeni, ogranicza zdolność miar podobieństwa do rozróżniania przykładów. Gao et al. wskazali, że degeneracja reprezentacji stanowi systematyczną konsekwencję treningu modeli generatywnych z funkcją straty opartą na wiarygodności \cite{gao2019representation}.

W odpowiedzi na to zjawisko opracowano metody przetwarzania końcowego przywracające izotropię przestrzeni. Metoda \textit{All-but-the-Top} usuwa z reprezentacji rzut na główne składowe odpowiadające za dominującą wariancję \cite{abtt}, natomiast wybielanie przekształca rozkład reprezentacji do postaci o zerowej średniej i jednostkowej macierzy kowariancji \cite{whiteningbert}. Odmienny wariant tej koncepcji stanowi \textit{concept whitening}, w którym transformacja wbudowana jest w architekturę modelu, dzięki czemu poszczególne osie przestrzeni odpowiadają zadanym pojęciom \cite{whitening}.

Analiza geometryczna znajduje zastosowanie również w kontekście modeli wielomodalnych i wielojęzycznych. Liang et al. opisali zjawisko \textit{modality gap}, polegające na tym, że reprezentacje różnych modalności w modelach uczonych kontrastywnie zajmują rozłączne obszary wspólnej przestrzeni \cite{liang2022mind}. Chang et al. wykazali natomiast, że w modelach wielojęzycznych wyodrębnić można podprzestrzenie specyficzne dla poszczególnych języków oraz podprzestrzeń wspólną, kodującą informację niezależną od języka \cite{chang2022geometry}. Do porównywania reprezentacji generowanych przez różne warstwy lub różne modele wykorzystywane są miary oparte na analizie korelacji kanonicznej, takie jak SVCCA \cite{raghu2017svcca}, oraz miara CKA \cite{pmlr-v97-kornblith19a}.

Metody opisane w tej części stanowią bezpośrednią podstawę metodologii ewaluacji przyjętej w rozdziale \ref{sec:latent}, w którym reprezentacje modeli ASR poddane zostały analizie pod kątem anizotropii, wpływu przetwarzania końcowego oraz separacji przykładów semantycznie podobnych i odległych.

\subsection{Podejścia probingowe i dekompozycja reprezentacji}

Metody sondujące (ang. \textit{probing}) opierają się na prostym rozumowaniu: jeżeli w reprezentacji wytworzonej przez daną warstwę modelu zakodowana jest pewna informacja, powinna ona dać się z niej odczytać za pomocą nieskomplikowanego klasyfikatora. Skuteczność takiego klasyfikatora, najczęściej liniowego, staje się wówczas miarą tego, na ile badana własność jest w reprezentacji dostępna \cite{alain2017probes}. Conneau et al. przenieśli tę ideę na poziom całych zdań, opracowując zestaw zadań pozwalających sprawdzić, które własności lingwistyczne -- od długości zdania po głębokość drzewa składniowego -- zachowują osadzenia zdaniowe \cite{conneau2018cram}. Belinkov i Glass zwracają jednak uwagę na zasadnicze ograniczenie tego podejścia: powodzenie sondy świadczy wyłącznie o tym, że informacja w reprezentacji się znajduje, nie zaś o tym, że model rzeczywiście się nią posługuje \cite{belinkov2019analysis}.

W domenie mowy najistotniejszym przykładem analizy probingowej jest praca Pasada et al., w której zbadano kolejne warstwy modelu Wav2Vec 2.0 \cite{pasad2021layer}. Wykazano, że model wykazuje charakterystykę zbliżoną do autoenkodera, przy czym warstwy środkowe kodują informację fonetyczną i leksykalną w stopniu najwyższym, natomiast warstwy końcowe wracają w kierunku reprezentacji zbliżonych do wejściowych. Obserwacja ta ma bezpośrednie znaczenie dla doboru warstwy wykorzystywanej jako źródło reprezentacji w architekturze proponowanej w niniejszej pracy. Pokrewny kierunek badań zakłada wymuszenie interpretowalności już na poziomie architektury, poprzez dyskretyzację reprezentacji do skończonego inwentarza jednostek, który można zestawiać z kategoriami fonetycznymi. Zhou et al. porównali stosowane w tym celu warianty kwantyzacji wektorowej, analizując ich wpływ na jakość reprezentacji mowy \cite{Zhou2020ACO}.

Podejścia jakościowe koncentrują się na rozkładzie reprezentacji na komponenty poddające się interpretacji. Punktem wyjścia jest obserwacja, że pojedyncze neurony sieci są zazwyczaj polisemantyczne, to znaczy aktywują się w odpowiedzi na wiele niepowiązanych ze sobą zjawisk, co wynika ze zjawiska superpozycji cech w przestrzeni o ograniczonej wymiarowości. Bricken et al. wykazali, że rzadkie autoenkodery (ang. \textit{sparse autoencoders}, SAE) trenowane na aktywacjach modelu językowego pozwalają na rozłożenie ich na słownik cech znacznie bliższych monosemantyczności \cite{bricken2023monosemanticity}. Wyniki te zostały potwierdzone i rozszerzone w pracy Hubena et al., w której wykazano wyższą interpretowalność cech uzyskanych metodą SAE w porównaniu z alternatywnymi metodami dekompozycji \cite{huben2024sparse}.

Podejście to znajduje coraz szersze zastosowanie w domenie audio. Paek et al. zastosowali rzadkie autoenkodery do przestrzeni cech ukrytych modeli audio, wyodrębniając składowe odpowiadające interpretowalnym zjawiskom akustycznym i muzycznym \cite{paek2025learninginterpretablefeaturesaudio}. Varshavsky-Hassid et al. wykazali natomiast obecność semantycznie nacechowanych kierunków w przestrzeni cech ukrytych dyfuzyjnych modeli syntezy mowy \cite{varshavskyhassid2024semanticlatentspacediffusionbased}. Prace te wskazują, że przestrzenie cech ukrytych modeli przetwarzających mowę zawierają strukturę wykraczającą poza poziom akustyczny, co stanowi bezpośrednią motywację dla badań prowadzonych w niniejszej pracy.

\subsection{Wyjaśnialność w systemach wyszukiwania opartych na mowie}

W architekturze kaskadowej transkrypcja stanowi pośredni artefakt umożliwiający kontrolę działania systemu przez użytkownika. Architektura pomijająca etap transkrypcji tej możliwości nie oferuje, co stanowi jej istotne ograniczenie. Analiza reprezentacji stanowi częściową rekompensatę tego ograniczenia. Pozwala ona na uzasadniony dobór modelu ASR i warstwy stanowiącej źródło reprezentacji, a także na diagnostykę własności przestrzeni wejściowej przed treningiem warstwy projekcji. Zestawienie metod omówionych w tej sekcji przedstawiono w Tabeli \ref{tab:xaiMetody}.

\begin{table}[htbp]
\centering
\small
\begin{tabular}{@{}L{4.0cm}L{6.6cm}L{3.2cm}@{}}
\toprule
\textbf{Metoda} & \textbf{Mierzona własność} & \textbf{Charakter} \\
\midrule
Klasyfikator probingowy & liniowa dostępność informacji w danej warstwie & ilościowy, nadzorowany \\
CKA, SVCCA & podobieństwo reprezentacji między warstwami i modelami & ilościowy \\
MMD na trypletach & separacja przykładów pozytywnych i negatywnych & ilościowy \\
Podobieństwo kosinusowe, efekt stożka & izotropia przestrzeni, dystans międzymodalny & ilościowy, geometryczny \\
ABTT, wybielanie & wpływ przetwarzania końcowego na zachowanie semantyki & ilościowy \\
Rzadkie autoenkodery (SAE) & interpretowalne składowe przestrzeni cech ukrytych & jakościowy \\
\bottomrule
\end{tabular}
\caption{Metody analizy reprezentacji wykorzystywane do oceny ich semantyczności.}
\label{tab:xaiMetody}
\end{table}


% ---------------------------------------------------
\section{Systemy wyszukiwania informacji w oparciu o sygnał audio}

Rosnące zapotrzebowanie na interfejsy głosowe motywuje rozwój systemów wyszukiwania informacji zdolnych do przetwarzania zapytań audio. Tradycyjne podejście opiera się na architekturze kaskadowej: sygnał mowy jest najpierw transkrybowany przez system ASR, a następnie uzyskany tekst wykorzystywany jest w standardowym procesie wyszukiwania. Rozwiązanie to jest jednak podatne na propagację błędów rozpoznawania mowy, szczególnie w przypadku terminologii specjalistycznej, różnych akcentów lub nagrań niskiej jakości.

W odpowiedzi na te ograniczenia powstały dwa główne nurty badawcze. Pierwszy koncentruje się na architekturach end-to-end, które uczą się bezpośredniego mapowania sygnału audio na reprezentacje semantyczne, omijając etap transkrypcji. Drugi wykorzystuje podejścia multimodalne, tworzące wspólną przestrzeń reprezentacji dla danych audio, tekstowych i obrazu.

\subsection{Podejścia end-to-end}

\textbf{SpeechDPR.} Model SpeechDPR (ang. \textit{Speech Dense Passage Retriever}) \cite{lin2024speechdprendtoendspokenpassage} stanowi pierwsze rozwiązanie typu end-to-end dla problemu wyszukiwania fragmentów mowy w zadaniach odpowiadania na pytania. Model uczy się semantycznych reprezentacji na poziomie zdaniowym poprzez destylację wiedzy z systemu kaskadowego składającego się z nienadzorowanego ASR oraz tekstowego retriever'a opartego na gęstych reprezentacjach. Proces ten nie wymaga ręcznie transkrybowanych danych mowy, co stanowi istotną zaletę w kontekście języków o ograniczonych zasobach. Schemat procesu treningu przedstawiono na Rysunku \ref{fig:speechdprArchitecture}.

\begin{figure}[htbp]
    \centering
    \begin{tikzpicture}[
    scale=0.6,
    transform shape,
    font=\small,
    block/.style={
        draw,
        rounded corners=2pt,
        thick,
        minimum width=2.8cm,
        minimum height=1.0cm,
        align=center
    },
    data/.style={block, fill=gray!12},
    model/.style={block, fill=blue!20},
    txt/.style={block, fill=purple!20},
    stud/.style={block, fill=teal!25},
    loss/.style={block, fill=yellow!25},
    arrow/.style={-{Stealth[]}, thick},
    phase/.style={draw=gray, dashed, rounded corners=4pt, inner sep=0.4cm}
]

% ============================================================
% WSPÓLNE WEJŚCIE
% ============================================================
\node[data] (src) {Audio zapytań\\i fragmentów};

% ============================================================
% NAUCZYCIEL: SYSTEM KASKADOWY
% ============================================================
\node[model, above right=1.2cm and 1.4cm of src] (asr) {Nienadzorowany\\model ASR};
\node[data, right=1.2cm of asr] (text) {Tekst};
\node[txt, right=1.2cm of text] (dpr) {Tekstowy retriever\\(DPR)};
\node[data, right=1.2cm of dpr] (scoret) {Oceny\\podobieństwa};

\draw[arrow] (src.east) -- ++(0.6,0) |- (asr.west);
\draw[arrow] (asr) -- (text);
\draw[arrow] (text) -- (dpr);
\draw[arrow] (dpr) -- (scoret);

% ============================================================
% UCZEŃ: SPEECHDPR
% ============================================================
\node[stud, below right=1.2cm and 1.4cm of src] (enc) {Enkoder mowy\\(SpeechDPR)};
\node[data, right=1.2cm of enc] (rep) {Reprezentacje\\zdaniowe};
\node[data, right=1.2cm of rep] (scores) {Podobieństwo\\kosinusowe};

\draw[arrow] (src.east) -- ++(0.6,0) |- (enc.west);
\draw[arrow] (enc) -- (rep);
\draw[arrow] (rep) -- (scores);

% ============================================================
% STRATA DESTYLACYJNA
% ============================================================
\node[loss, right=1.4cm of scoret.east |- src] (loss) {Strata\\destylacyjna};

\draw[arrow] (scoret.east) -| (loss.north);
\draw[arrow] (scores.east) -| (loss.south);

% ============================================================
% RAMKI
% ============================================================
\node[phase, fit=(asr)(scoret),
      label={[font=\small\itshape, gray]above:Nauczyciel -- system kaskadowy (bez ręcznych transkrypcji)}] {};
\node[phase, fit=(enc)(scores),
      label={[font=\small\itshape, gray]below:Uczeń -- model end-to-end}] {};

\end{tikzpicture}

    \caption{Schemat treningu modelu SpeechDPR. Kaskadowy nauczyciel, złożony z nienadzorowanego modelu ASR i tekstowego retrievera, dostarcza ocen podobieństwa stanowiących cel destylacji dla enkodera mowy uczonego bezpośrednio na sygnale audio.}
    \label{fig:speechdprArchitecture}
\end{figure}

\textbf{SpeechRAG.} SpeechRAG \cite{min2025speechretrievalaugmentedgenerationautomatic} rozszerza architekturę RAG o bezpośrednie przetwarzanie audio. System składa się z retriever'a audio oraz generatora opartego na modelu języka mowy (SLM, ang. \textit{Speech Language Model}). Kluczową innowacją jest dostrojenie pretrenowanego enkodera mowy do roli adaptera podłączonego do zamrożonego retriever'a tekstowego. Dzięki wyrównaniu przestrzeni embeddingów retriever audio może wyszukiwać fragmenty na podstawie zapytań tekstowych. W fazie generacji model SLM otrzymuje jako kontekst oryginalne fragmenty audio, eliminując straty informacyjne wynikające z konwersji mowa-tekst. Architekturę systemu przedstawiono na Rysunku \ref{fig:speechragArchitecture}.

\begin{figure}[htbp]
    \centering
    \begin{tikzpicture}[
    scale=0.6,
    transform shape,
    font=\small,
    block/.style={
        draw,
        rounded corners=2pt,
        thick,
        minimum width=2.8cm,
        minimum height=1.0cm,
        align=center
    },
    data/.style={block, fill=gray!12},
    adapt/.style={block, fill=teal!25},
    frozen/.style={block, fill=blue!20},
    store/.style={block, fill=orange!25},
    ann/.style={block, fill=yellow!25},
    gen/.style={block, fill=green!20},
    arrow/.style={-{Stealth[]}, thick},
    phase/.style={draw=gray, dashed, rounded corners=4pt, inner sep=0.4cm}
]

% ============================================================
% FAZA WYSZUKIWANIA
% ============================================================
\node[data] (query) {Zapytanie\\tekstowe};
\node[frozen, right=1.8cm of query] (enct) {Enkoder tekstu\\(zamrożony)};
\node[data, right=1.8cm of enct] (qvec) {Wektor\\zapytania};
\node[ann, right=1.8cm of qvec] (search) {Wyszukiwanie ANN};

\draw[arrow] (query) -- (enct);
\draw[arrow] (enct) -- (qvec);
\draw[arrow] (qvec) -- (search);

% ============================================================
% FAZA INDEKSOWANIA
% ============================================================
\node[data, above=2.6cm of query] (chunks) {Fragmenty\\audio};
\node[adapt, above=2.6cm of enct] (adapter) {Enkoder mowy --\\adapter (trenowany)};
\node[store, above=2.6cm of search] (db) {Wektorowa\\baza danych};

\draw[arrow] (chunks) -- (adapter);
\draw[arrow] (adapter) -- (db);
\draw[arrow] (db) -- node[right, font=\scriptsize] {indeks} (search);

% wyrównanie przestrzeni embeddingów
\draw[{Stealth[]}-{Stealth[]}, dotted, thick, gray] (enct.north) --
    node[left, pos=0.78, font=\scriptsize\itshape, text=gray]
    {wspólna przestrzeń} (adapter.south);

% ============================================================
% FAZA GENERACJI
% ============================================================
\node[data, below=2.6cm of query] (topk) {Top-$k$ fragmentów\\audio};
\node[gen, right=2.6cm of topk] (slm) {Model języka mowy\\(SLM)};
\node[data, right=1.8cm of slm] (ans) {Odpowiedź};

\draw[arrow] (search.south) -- ++(0,-0.8) -| (topk.north);
\draw[arrow] (topk) -- node[above, font=\scriptsize] {kontekst audio} (slm);
\draw[arrow] (slm) -- (ans);

% ============================================================
% RAMKI FAZ
% ============================================================
\node[phase, fit=(chunks)(db),
      label={[font=\small\itshape, gray]above:Faza indeksowania}] {};
\node[phase, fit=(query)(search),
      label={[font=\small\itshape, gray]above left:Faza wyszukiwania}] {};
\node[phase, fit=(topk)(ans),
      label={[font=\small\itshape, gray]below left:Faza generacji}] {};

\end{tikzpicture}

    \caption{Architektura systemu SpeechRAG. Enkoder mowy pełniący rolę adaptera jest dostrajany tak, aby fragmenty audio trafiały do przestrzeni wektorowej zamrożonego enkodera tekstu. Wyszukane fragmenty audio stanowią bezpośredni kontekst modelu języka mowy.}
    \label{fig:speechragArchitecture}
\end{figure}

\subsection{Podejścia multimodalne}

Budowa systemów wyszukiwania opartych na wspólnej przestrzeni reprezentacji dla audio i tekstu jest ułatwiona dzięki modelom takim jak Data2Vec (opisany w podrozdziale \ref{sec:modeleASR}), które oferują zunifikowane podejście do uczenia reprezentacji różnych modalności.

\textbf{WavRAG.} WavRAG \cite{wavRAG} umożliwia reprezentowanie sygnałów audio w tej samej przestrzeni semantycznej co dane tekstowe. Pozwala to na pełną integrację informacji mówionej i pisanej w ramach jednego systemu. Zapytania głosowe są kodowane jako wektory w przestrzeni cech ukrytych, a wyszukiwanie opiera się na bezpośrednim dopasowywaniu embeddingów akustycznych i tekstowych za pomocą podobieństwa kosinusowego. Architekturę systemu przedstawiono na Rysunku \ref{fig:wavragArchitecture}.

\begin{figure}[htbp]
    \centering
    \begin{tikzpicture}[
    scale=0.6,
    transform shape,
    font=\small,
    block/.style={
        draw,
        rounded corners=2pt,
        thick,
        minimum width=2.8cm,
        minimum height=1.0cm,
        align=center
    },
    data/.style={block, fill=gray!12},
    embm/.style={block, fill=purple!20},
    store/.style={block, fill=orange!25},
    ann/.style={block, fill=yellow!25},
    gen/.style={block, fill=green!20},
    arrow/.style={-{Stealth[]}, thick},
    phase/.style={draw=gray, dashed, rounded corners=4pt, inner sep=0.4cm}
]

% ============================================================
% FAZA WYSZUKIWANIA
% ============================================================
\node[data] (query) {Zapytanie\\głosowe};
\node[embm, right=1.8cm of query] (encq) {Enkoder\\multimodalny};
\node[data, right=1.8cm of encq] (qvec) {Wektor\\zapytania};
\node[ann, right=1.8cm of qvec] (search) {Wyszukiwanie ANN\\(podobieństwo kosinusowe)};

\draw[arrow] (query) -- (encq);
\draw[arrow] (encq) -- (qvec);
\draw[arrow] (qvec) -- (search);

% ============================================================
% FAZA INDEKSOWANIA
% ============================================================
\node[data, above=4.4cm of query] (docs) {Dokumenty\\tekstowe};
\node[data, below=0.7cm of docs] (recs) {Nagrania\\audio};
\node[embm, above=3.6cm of encq] (enci) {Enkoder\\multimodalny};
\node[store, above=3.6cm of search] (db) {Wspólna przestrzeń\\wektorowa};

\draw[arrow] (docs.east) -- ++(0.5,0) |- (enci.west);
\draw[arrow] (recs.east) -- ++(0.5,0) |- (enci.west);
\draw[arrow] (enci) -- (db);
\draw[arrow] (db) -- node[right, font=\scriptsize] {indeks} (search);

% wspólne wagi enkodera dla obu faz
\draw[{Stealth[]}-{Stealth[]}, dotted, thick, gray] (encq.north) --
    node[left, pos=0.4, font=\scriptsize\itshape, text=gray] {wspólne wagi} (enci.south);

% ============================================================
% FAZA GENERACJI
% ============================================================
\node[data, below=2.6cm of query] (topk) {Fragmenty tekstowe\\i audio};
\node[gen, right=2.6cm of topk] (gen) {Model\\generacyjny};
\node[data, right=1.8cm of gen] (ans) {Odpowiedź};

\draw[arrow] (search.south) -- ++(0,-0.8) -| (topk.north);
\draw[arrow] (topk) -- node[above, font=\scriptsize] {kontekst mieszany} (gen);
\draw[arrow] (gen) -- (ans);

% ============================================================
% RAMKI FAZ
% ============================================================
\node[phase, fit=(docs)(recs)(enci)(db),
      label={[font=\small\itshape, gray]above:Faza indeksowania}] {};
\node[phase, fit=(query)(search),
      label={[font=\small\itshape, gray]above left:Faza wyszukiwania}] {};
\node[phase, fit=(topk)(ans),
      label={[font=\small\itshape, gray]below left:Faza generacji}] {};

\end{tikzpicture}

    \caption{Architektura systemu WavRAG. Ten sam enkoder multimodalny obsługuje obie modalności, dzięki czemu dokumenty tekstowe i nagrania audio trafiają do wspólnej przestrzeni wektorowej, a wyszukane fragmenty obu typów stanowią kontekst modelu generacyjnego.}
    \label{fig:wavragArchitecture}
\end{figure}

\textbf{Wyszukiwanie multimodalne dla ASR.} Kolehmainen et al. \cite{kolehmainen2024multimodalretrievallargelanguage} zbadali zastosowanie wyszukiwania multimodalnego do poprawy jakości rozpoznawania mowy. Zaproponowany system wykorzystuje dwie techniki: kNN-LM (ang. \textit{k-Nearest Neighbors Language Model}) oraz mechanizmy cross-attention. Wyszukiwanie oparte na mowie prowadzi do redukcji WER o nawet 50\% w porównaniu z bazowym modelem. System osiągnął wyniki state-of-the-art w zadaniu rozpoznawania mowy na zbiorze Spoken-Squad. Sposób włączenia wyszukanych fragmentów do procesu rozpoznawania przedstawiono na Rysunku \ref{fig:multimodalAsrArchitecture}.

\begin{figure}[htbp]
    \centering
    \begin{tikzpicture}[
    scale=0.6,
    transform shape,
    font=\small,
    block/.style={
        draw,
        rounded corners=2pt,
        thick,
        minimum width=2.8cm,
        minimum height=1.0cm,
        align=center
    },
    data/.style={block, fill=gray!12},
    model/.style={block, fill=blue!20},
    dec/.style={block, fill=purple!20},
    ret/.style={block, fill=yellow!25},
    store/.style={block, fill=orange!25},
    fuse/.style={block, fill=teal!25},
    arrow/.style={-{Stealth[]}, thick},
    phase/.style={draw=gray, dashed, rounded corners=4pt, inner sep=0.4cm}
]

% ============================================================
% ŚCIEŻKA ROZPOZNAWANIA MOWY
% ============================================================
\node[data] (audio) {Nagranie\\audio};
\node[model, right=1.4cm of audio] (enc) {Enkoder ASR};
\node[dec, right=1.4cm of enc] (dec) {Dekoder ASR};
\node[fuse, right=1.4cm of dec] (interp) {Interpolacja\\rozkładów};
\node[data, right=1.4cm of interp] (text) {Transkrypcja};

\draw[arrow] (audio) -- (enc);
\draw[arrow] (enc) -- (dec);
\draw[arrow] (dec) -- node[above, font=\scriptsize] {$p_{\mathrm{LM}}$} (interp);
\draw[arrow] (interp) -- (text);

% ============================================================
% ŚCIEŻKA WYSZUKIWANIA MULTIMODALNEGO
% ============================================================
\node[store, below=2.8cm of audio] (kb) {Baza wiedzy\\(audio i tekst)};
\node[ret, below=2.8cm of enc] (ret) {Wyszukiwanie\\multimodalne};
\node[data, below=2.8cm of dec] (topk) {Top-$k$\\fragmentów};
\node[data, below=2.8cm of interp] (knn) {Rozkład\\$k$-NN};

\draw[arrow] (kb) -- (ret);
\draw[arrow] (ret) -- (topk);
\draw[arrow] (topk) -- node[above, font=\scriptsize] {$p_{k\mathrm{NN}}$} (knn);

\draw[arrow] (audio.south) -- ++(0,-1.4)
    -| node[pos=0.25, above, font=\scriptsize] {zapytanie} (ret.north);

% ============================================================
% MECHANIZMY INTEGRACJI
% ============================================================
\draw[arrow, thick] (topk.north) --
    node[right, font=\scriptsize\itshape] {\textit{cross-attention}} (dec.south);
\draw[arrow, thick] (knn.north) --
    node[right, font=\scriptsize\itshape] {kNN-LM} (interp.south);

\node[phase, fit=(kb)(knn),
      label={[font=\small\itshape, gray]below:Wyszukiwanie multimodalne}] {};

\end{tikzpicture}

    \caption{Wyszukiwanie multimodalne wspierające rozpoznawanie mowy. Wyszukane fragmenty oddziałują na model ASR dwiema drogami: poprzez mechanizm uwagi krzyżowej w dekoderze oraz poprzez interpolację rozkładu prawdopodobieństwa tokenów z rozkładem wyznaczonym metodą $k$-NN.}
    \label{fig:multimodalAsrArchitecture}
\end{figure}

\subsection{Charakterystyka i porównanie podejść}

Przedstawione podejścia łączy dążenie do eliminacji lub ograniczenia roli tradycyjnego ASR jako pośredniego etapu przetwarzania. Zarówno architektury end-to-end, jak i multimodalne wykazują przewagę nad systemami kaskadowymi w warunkach niskiej jakości rozpoznawania mowy. Bezpośrednie uczenie reprezentacji semantycznych z sygnału audio przynosi dwie kluczowe korzyści: redukuje propagację błędów transkrypcji oraz otwiera możliwość zachowania informacji prozodycznych i paralingwistycznych, traconych w procesie konwersji mowa-tekst. Żadna z omówionych prac nie wykazuje jednak, by informacja prozodyczna była w tych systemach faktycznie wykorzystywana przy wyszukiwaniu. Stanowi to fundamentalną zmianę paradygmatu w projektowaniu systemów wyszukiwania informacji opartych na danych mowy.
